В ходе обучения на протяжении всей своей жизни были освоены или находятся в процессе освоения теоретические знания как из области цифровых технологий, так и из области прикладных задач, в числе которых высшая математика, электротехника, общая история России, механика и др.

Тем не менее, особая часть знаний сконцентирована вокруг цифровых технологий, на что делается особый упор в свете настоящих событий. Так, благодаря имеющемуся багажу знаний, были реализованы слудеющие проекты:
\begin{itemize}
    \item Информационный справочник на основе открытого gitbook для организации информации о процессе глубокой модификации Android устройств в виде установки сторонней прошивки, смены ядра, открытия root доступа и др. Данный проект был разработан по личной инициативе в рамках открытого специализированного форума.
    \item Информационный справочник на основе веб-сайта для организации школьных материалов, свободно размещенный в сети Интернет. Данный проект был разработан в рамках школьных проектов за 10--11 классы.
    \item Трехмерное моделирование и оформление конструкторской документации изделий в системе Компас-3D. Данный проект был реализован в рамках курсовой работы по дисциплине <<инженерная графика>>.
    \item Система управления базой данных на языке программирования C для организации картотеки чертежей с псевдографическим интерфейсом. Данный проект был разработан в рамках курсовой работы по дисциплине <<информационные технологии>>.
    \item Интерфейс для работы с библиотекой BLAS (Fortran) в C++. Данный проект был разработан по личной инициативе в образовательных целях.
    \item Алгоритм поиска кратчайшего пути по Форду-Беллману на MatLab. Данный проект был реализован в рамках курсовой работы по дисциплине <<программирование и основы алгоритмизации>>.
    \item Среда формирования типизированных отчетов в соответствии с ГОСТ 7.32-2017 (об отчета о НИР) и требованиями оформления СПбГЭТУ <<ЛЭТИ>>. Данный проект разрабатывается по личной инициативе в целях упрощения составления отчетов.
\end{itemize}

Анализ личностных особенностей на одной из практик настоящей дисциплины показывает следующие результаты:
\begin{itemize}
    \item \textit{Я думаю, что я западный человек, потому что я всегда планирую. Планы и расписание кажутся мне удобными инструментами, не дающими забыть или пропустить важные вещи. Если не думать о будущем, то на него нельзя будет повлиять как мне хочется. И всё же, действительно, иногда можно дать себе слабину и кайфовать в моменте.}
    \item \textit{Я думаю, что мне сложно отнести себя либо к жаворонку, либо сове. Я могу быть одинаково продуктивным как утром, так и вечером по необходимости. Были и есть случаи, когда я хорошо работал как ранним утром, так и вечером (хотя я стараюсь заканчивать к 00:00; не допускаю работу после 01:00, чтобы не сбивать режим).}
    \item \textit{Я думаю, что я определенно визуал. Про кинестетика вообще сложно говорить, поскольку формулы не пощупать. Если же кинестетика — это про письмо ручкой по бумаге, то определенно да. В таком случае «визуал» и «кинестетик» обычно дополняют друг друга. Аудиал крайне мало присутствует в моей жизни, разве что мнемоники могут хорошо работать.}
    \item \textit{Я думаю, что мне сложно отнести себя к какому-то типу. Конечно, заполнять однотипные формы — ужасно, но ведь можно автоматизировать. При этом вводить что-то новое на замену старого без практической необходимости не имеет смысла. В то же самое время использование устаревших систем так же плохо. Необходим правильный баланс в сангвинике-флегматике.}
    \item \textit{[e/i]. Я думаю, что я экстраверт. Я действительно «просто говорю». Между офисом и удаленкой я бы выбрал офис, поскольку сидеть одному всё время дома одиноко и скучно.}
    \item \textit{[s/n]. Я думаю, что больше сенсорик, однако между жанрами повседневность и фэнтези я бы выбрал второе, поскольку зачем мне читать/смотреть на ещё одну такую же реальность? Буквальность для меня предпочтительнее намеков, но играться словами тоже весело.}
    \item \textit{[t/f]. Я думаю, что я больше feeler. Нельзя управлять командой, не зная её.}
    \item \textit{[j/p]. Я думаю, что я judger. Действительно, спланировав будущее, можно его спрогнозировать. Я действительно с великой точностью знаю, какой предмет где находится.}
\end{itemize}

В то же самое время, ранее пройденнный тест на веб-сайте \url{16personalities.com} показывает результат в виде ENTJ-A. Однако в связи с наличием более полного анализа, проведенного в рамках практики настоящей дисциплины, отбросим данный результат и сфокусируемся на самоанализе выше:
\begin{itemize}
    \item \textit{Характер.} Последовательный, прямолинейный. Деятельность базируется на жестком планировании и соблюдении графиков, что минимизирует риски пропуска критических задач.
    \item \textit{Принятие решений.} Основано на учтении текущей загруженности и дальнейший планов. Не импульсивное, но расчетливое.
    \item \textit{Отношение к людям.} Нет никаких ярко выраженного отрицательного отношения. Нейтральное отношение ко всем, кто не представляет интереса или не вмешивается в мою зону комфорта.
    \item \textit{Способности к коммуникации.} Хорошие способности к коммуникации, нет внутренних ограничителей, мешающих строить общение с людьми.
\end{itemize}

Составим таблицу на основе данных из личностного самоанализа. См. \cref{tab:swot_analysis}.

\begin{longtblr}[
    caption = {Анализ сильных и слабых сторон личности, возможностей и угроз},
    label = {tab:swot_analysis},
]{
    colspec = {Q[c,m] X[j,t] X[j,t]},
    hlines, vlines,
    row{1,3} = {bg=gray9},
    rowhead = 1,
}
    \SetCell[r=2]{m} \rotatebox{90}{Внутренняя среда} & \SetCell{c} Сильные стороны личности & \SetCell{c} Слабые стороны личности \\
    & Навыки программирования и проектирования; развитое системное планирование; коммуникабельность; высокая адаптивность к нагрузкам. & 
    Низкая толерантность к монотонным, неавтоматизированным процессам; зависимость продуктивности от наличия четкого плана. \\
    \SetCell[r=2]{m} \rotatebox{90}{Внешняя среда} & \SetCell{c} Возможности & \SetCell{c} Угрозы \\
    & Востребованность специалистов по автоматизации систем управления; государственная поддержка IT-отрасли; доступ к информационным ресурсам (в т.ч. университета). & 
    Конкуренция со стороны опытных специалистов; рыночная нестабильность; хаотичная политика государства. \\
\end{longtblr}

На пересечении выявленных характеристик сформулированы два варианта развития:
\begin{enumerate}
    \item \textit{Основной вариант:} Трудоустройство в инженерную компанию на должность инженера-проектировщика или разработчика систем автоматизации.
    \item \textit{Альтернативный вариант:} Самозанятость в сфере разработки программного обеспечения и веб-технологий.
\end{enumerate}
